\documentclass[10pt,xcolor=dvipsnames]{beamer}
\usetheme{Madrid}
\usecolortheme[RGB={0,82,155}]{structure}
\setbeamertemplate{navigation symbols}{}

\usepackage[utf8]{inputenc}
\usepackage[spanish]{babel}
\usepackage{amsmath,amssymb}
\usepackage{hyperref}
\usepackage{tikz}
\usetikzlibrary{arrows.meta}
\usepackage{booktabs}

\title[Sistemas Numéricos]{Sistemas Numéricos en la Historia}
\subtitle{Historia de la Matemática --- MM-517}
\author{}
\institute[UNAH]{Universidad Nacional Autónoma de Honduras \\ Facultad de Ciencias \\ Departamento de Matemáticas}
\date{III Período, 2026}

\AtBeginSection[]{%
  \begin{frame}<beamer>\frametitle{Contenido}
    \tableofcontents[currentsection]
  \end{frame}
}

\begin{document}

\begin{frame}
  \titlepage
\end{frame}

\begin{frame}{Contenido}
  \tableofcontents
\end{frame}

%% ======================================================================
\section{Sistemas numéricos en la prehistoria}
%% ======================================================================

\begin{frame}{Paleolítico Superior: el origen del conteo}
  \begin{itemize}
    \item Período: \textbf{Paleolítico Superior} (\textasciitilde 50{,}000--10{,}000 BCE)
    \item Notación más antigua: \textbf{muescas en hueso} (\textit{tally marks})
    \item \textbf{Hueso de Lebombo}: \textasciitilde 44{,}000 años
      \begin{itemize}
        \item Border Cave (Sudáfrica)
        \item 29 muescas: posible calendario lunar
      \end{itemize}
    \item \textbf{Hueso de Ishango}: \textasciitilde 20{,}000 años
      \begin{itemize}
        \item Fíbula de babuino, RD Congo (descubierto 1960)
        \item Tres columnas de muescas: posible aritmética
      \end{itemize}
    \item Conteo corporal: bases 5 y 10
  \end{itemize}
\end{frame}

\begin{frame}{Prehistoria: artefactos clave}
  \vspace{1cm}
  \begin{center}
    \fbox{\begin{minipage}[c][5cm][c]{0.85\textwidth}
      \centering \textit{( Espacio para imágenes: \\ hueso de Ishango, hueso de Lebombo, \\ marcas en cuevas como Lascaux )}
    \end{minipage}}
  \end{center}
\end{frame}

%% ======================================================================
\section{El período neolítico}
%% ======================================================================

\begin{frame}{Neolítico y casos clave (Ötzi, Egtved)}
  \begin{itemize}
    \item \textbf{Período Neolítico}: \textasciitilde 10{,}000--4{,}500 BCE (Cercano Oriente)
    \item \textbf{Ötzi el Hombre de Hielo}: \textasciitilde 3300 BCE (Edad del Cobre / Neolítico tardío)
      \begin{itemize}
        \item Hallado en los Alpes de Ötztal (Austria/Italia), 1991
        \item Hacha de cobre + herramientas de piedra
        \item Más antiguo que las pirámides egipcias y Stonehenge
      \end{itemize}
    \item \textbf{Niña de Egtved}: \textasciitilde 1370 BCE --- \textit{Edad del Bronce}
      \begin{itemize}
        \item Hallada en Dinamarca (1921); momificación en turbera
        \item Edad estimada: 16--18 años
      \end{itemize}
  \end{itemize}
\end{frame}

\begin{frame}{Imágenes y nota}
  \vspace{0.5cm}
  \begin{center}
    \fbox{\begin{minipage}[c][4.5cm][c]{0.85\textwidth}
      \centering \textit{( Espacio para imágenes: \\ reconstrucción de Ötzi, sarcófago de la Niña de Egtved )}
    \end{minipage}}
  \end{center}
  \vspace{0.3cm}
  \begin{alertblock}{Nota}
    La Niña de Egtved pertenece a la \textbf{Edad del Bronce}, no al Neolítico estricto.
  \end{alertblock}
\end{frame}

%% ======================================================================
\section{Matemáticas del Neolítico}
%% ======================================================================

\begin{frame}{¿Qué matemáticas se hacían?}
  \begin{itemize}
    \item \textbf{Sin escritura} ni notación simbólica formal
    \item Herramientas principales:
      \begin{itemize}
        \item Cuerdas con nudos
        \item Muescas en palos y huesos (\textit{tally sticks})
        \item Marcas en cuevas y megalitos
      \end{itemize}
    \item \textbf{Observación astronómica}:
      \begin{itemize}
        \item Ciclos lunares (29 días)
        \item Equinoccios y solsticios
      \end{itemize}
    \item \textbf{Geometría práctica}:
      \begin{itemize}
        \item Stonehenge (\textasciitilde 3000--2000 BCE): alineamientos a solsticios; posible calendario solar de 365.25 días
        \item Triángulos pitagóricos (3-4-5, 5-12-13) en megalitos
      \end{itemize}
  \end{itemize}
\end{frame}

\begin{frame}{Artefactos y geometría}
  \vspace{1cm}
  \begin{center}
    \fbox{\begin{minipage}[c][5cm][c]{0.85\textwidth}
      \centering \textit{( Espacio para imágenes: \\ Stonehenge, tally stick neolítico, alineamientos )}
    \end{minipage}}
  \end{center}
\end{frame}

%% ======================================================================
\section{Sistema posicional babilónico}
%% ======================================================================

\begin{frame}{Sexagesimal y cuneiforme}
  \begin{itemize}
    \item \textbf{Origen}: Sumerios (\textasciitilde 3500--3000 BCE), Mesopotamia
    \item \textbf{Sistema sexagesimal} (base 60)
    \item \textbf{Escritura cuneiforme} sobre tablillas de arcilla
    \item \textbf{Posicional}: el valor depende de la posición
    \item Sub-base 10: dígitos 1--9 (cuñas) y 10 (esquina)
    \item Sin cero (espacio en blanco)
    \item \textbf{Períodos}:
      \begin{itemize}
        \item Sumerio: 3500--2000 BCE
        \item Babilónico antiguo: 2000--1595 BCE
        \item Babilónico nuevo: 1595--1155 BCE
        \item Seleúcida: 312--63 BCE
      \end{itemize}
  \end{itemize}
\end{frame}

\begin{frame}{Tablillas famosas y legado moderno}
  \begin{itemize}
    \item \textbf{Plimpton 322} (\textasciitilde 1900--1600 BCE): triadas pitagóricas
    \item Tablillas de Senkerah: cuadrados hasta 59, cubos hasta 32
    \item Tablas de recíprocos, raíces cuadradas, ecuaciones
    \item \textbf{Legado moderno}:
      \begin{itemize}
        \item 60 segundos/minuto, 60 minutos/hora
        \item 360$^{\circ}$ por circunferencia
        \item Coordenadas geográficas (latitud/longitud)
      \end{itemize}
  \end{itemize}
  \vspace{0.3cm}
  \begin{center}
    \fbox{\begin{minipage}[c][2cm][c]{0.55\textwidth}
      \centering \textit{( Imagen: Plimpton 322 )}
    \end{minipage}}
  \end{center}
\end{frame}

%% ======================================================================
\section{Numeración egipcia}
%% ======================================================================

\begin{frame}{Sistema jeroglífico}
  \begin{itemize}
    \item \textbf{Período}: \textasciitilde 3000 BCE -- 300 CE
    \item Sistema \textbf{decimal} (base 10)
    \item \textbf{Sin notación posicional}
    \item \textbf{Sin cero} (nfr $\neq$ cero matemático)
    \item \textbf{Símbolos jeroglíficos}:
      \begin{itemize}
        \item 1 = trazo vertical
        \item 10 = traba de ganado
        \item 100 = espiral de cuerda
        \item 1{,}000 = flor de loto
        \item 10{,}000 = dedo
        \item 100{,}000 = renacuajo
        \item 1{,}000{,}000 = dios Heh
      \end{itemize}
    \item 9{,}999 necesitaba 36 símbolos (suma, no posición)
  \end{itemize}
\end{frame}

\begin{frame}{Numeración hierática}
  \begin{itemize}
    \item Origen: \textasciitilde 3000 BCE; formalizada \textasciitilde 1800 BCE
    \item Usada en papiros: \textbf{Rhind} (\textasciitilde 1650 BCE), \textbf{Moscú} (\textasciitilde 1850 BCE)
    \item Versión cursiva de los jeroglíficos
    \item Coexistió con jeroglíficos por \textasciitilde 2000 años
    \item \textbf{Compactación}:
      \begin{itemize}
        \item 9{,}999 en jeroglífico: 36 símbolos
        \item 9{,}999 en hierático: 4 símbolos
      \end{itemize}
    \item Seis formas históricas distintas (evolución)
  \end{itemize}
  \vspace{0.3cm}
  \begin{center}
    \fbox{\begin{minipage}[c][2cm][c]{0.55\textwidth}
      \centering \textit{( Imagen: comparación jeroglífico vs hierático )}
    \end{minipage}}
  \end{center}
\end{frame}

%% ======================================================================
\section{Sistemas numéricos de América}
%% ======================================================================

\begin{frame}{Vigesimal y orígenes olmecas}
  \begin{itemize}
    \item \textbf{Origen}: cultura olmeca (\textasciitilde 1500--400 BCE), costa del Golfo de México
    \item \textbf{Base común}: vigesimal (base 20)
      \begin{itemize}
        \item Conteo con dedos de manos y pies
      \end{itemize}
    \item \textbf{Símbolos básicos}:
      \begin{itemize}
        \item Punto = 1
        \item Barra = 5
        \item Concha = 0
      \end{itemize}
    \item \textbf{Culturas}:
      \begin{itemize}
        \item Olmeca (madre), Zapoteca (Monte Albán)
        \item Maya, Azteca, Mixteca
      \end{itemize}
    \item Calendario de Cuenta Larga: usado desde \textasciitilde 36 BCE
  \end{itemize}
\end{frame}

\begin{frame}{Mapa cronológico}
  \vspace{1cm}
  \begin{center}
    \fbox{\begin{minipage}[c][5cm][c]{0.85\textwidth}
      \centering \textit{( Espacio para: \\ mapa de Mesoamérica con civilizaciones, \\ línea de tiempo de culturas )}
    \end{minipage}}
  \end{center}
\end{frame}

%% ======================================================================
\section{Sistema numérico y jeroglífico maya}
%% ======================================================================

\begin{frame}{Vigesimal posicional}
  \begin{itemize}
    \item \textbf{Período Clásico}: \textasciitilde 250--900 CE (sur de México, Guatemala, Belice)
    \item \textbf{Vigesimal} (base 20) posicional
    \item \textbf{Tres símbolos}:
      \begin{itemize}
        \item Concha (o glifo de pétalo) = 0
        \item Punto = 1
        \item Barra = 5
      \end{itemize}
    \item \textbf{Modificación}: 2$^{\circ}$ nivel es $18 \times 20 = 360$ (no $20 \times 20 = 400$), alineado con el año solar
    \item \textbf{Unidades temporales}:
      \begin{itemize}
        \item \textit{kin} = 1 día
        \item \textit{uinal} = 20 kin
        \item \textit{tun} = 18 uinal = 360 días
        \item \textit{katun} = 20 tun = 7{,}200 días
        \item \textit{baktun} = 20 katun = 144{,}000 días
      \end{itemize}
  \end{itemize}
\end{frame}

\begin{frame}{Calendarios y cero}
  \begin{itemize}
    \item \textbf{Tres calendarios}:
      \begin{itemize}
        \item \textit{Haab} (365 días, año vago)
        \item \textit{Tzolkin} (260 días, año sagrado)
        \item \textit{Cuenta Larga} (cronología absoluta)
      \end{itemize}
    \item \textbf{Cero matemático} usado desde \textasciitilde siglo III CE
    \item Primera fecha Cuenta Larga: \textbf{36 BCE} (Estela 2, Chiapa de Corzo)
    \item Notación vertical u horizontal
  \end{itemize}
  \vspace{0.3cm}
  \begin{center}
    \fbox{\begin{minipage}[c][2cm][c]{0.55\textwidth}
      \centering \textit{( Imagen: estela maya con glifos numéricos )}
    \end{minipage}}
  \end{center}
\end{frame}

%% ======================================================================
\section{Sistema azteca}
%% ======================================================================

\begin{frame}{Pictográfico vigesimal}
  \begin{itemize}
    \item \textbf{Imperio Mexica/Azteca (Triple Alianza)}: 1428--1521 CE
    \item \textbf{Vigesimal} (base 20) pictográfico
    \item \textbf{Símbolos}:
      \begin{itemize}
        \item Punto = 1
        \item Bandera / estandarte = 20
        \item Pluma = 400
        \item Bulto de copal = 8{,}000
      \end{itemize}
    \item \textbf{Uso}: tributos y registro comercial
    \item \textbf{Sin cero} explícito (a diferencia de Maya)
    \item Sistema \textbf{no posicional} estricto (suma)
  \end{itemize}
\end{frame}

\begin{frame}{Fuentes y comparación}
  \vspace{0.8cm}
  \begin{center}
    \fbox{\begin{minipage}[c][4.5cm][c]{0.85\textwidth}
      \centering \textit{( Espacio para: \\ Códice Mendoza, glifos numéricos aztecas, \\ comparación con Maya )}
    \end{minipage}}
  \end{center}
  \vspace{0.2cm}
  \begin{alertblock}{Diferencia con Maya}
    Los aztecas \textbf{sumaban}; los mayas tenían notación \textbf{posicional} y usaban el \textbf{cero}.
  \end{alertblock}
\end{frame}

%% ======================================================================
\section{Los quipus incas}
%% ======================================================================

\begin{frame}{Estructura y materiales}
  \begin{itemize}
    \item \textbf{Imperio Inca}: \textasciitilde 1438--1532 CE
    \item Uso más antiguo: \textasciitilde 2600 BCE (culturas andinas pre-Inca)
    \item \textbf{Cuerdas anudadas} (algodón o fibra de camélido: alpaca, llama, vicuña)
    \item \textbf{Información codificada}:
      \begin{itemize}
        \item Posición del nudo (decimal posicional)
        \item Color del cordón
        \item Tipo de nudo (simple, largo, en figura)
        \item Distancia entre nudos
        \item Orden de los cordones
      \end{itemize}
  \end{itemize}
\end{frame}

\begin{frame}{Uso, especialistas y escala}
  \begin{itemize}
    \item \textbf{Función principal}: numérica (decimal posicional)
    \item \textbf{Registros}: censos, tributos, calendario, producción, militar
    \item \textbf{Administradores}: \textit{quipucamayocs} (guardianes de la memoria)
    \item Estudios recientes sugieren posible contenido narrativo en algunos quipus
    \item \textbf{Escala}: desde pocos cordones hasta miles
  \end{itemize}
  \vspace{0.3cm}
  \begin{center}
    \fbox{\begin{minipage}[c][2cm][c]{0.55\textwidth}
      \centering \textit{( Imagen: quipu reconstituido, ejemplo de nudos )}
    \end{minipage}}
  \end{center}
\end{frame}

\end{document}
